\documentclass[11pt]{article}
\usepackage[margin=1in]{geometry}
\usepackage{hyperref}
\usepackage{booktabs}
\usepackage{listings}
\usepackage{xcolor}
\usepackage{amsmath}

\hypersetup{colorlinks=true, linkcolor=teal, urlcolor=teal, citecolor=teal}

\lstset{
  basicstyle=\ttfamily\small,
  breaklines=true,
  frame=single,
  rulecolor=\color{gray!40},
  backgroundcolor=\color{gray!5},
  columns=fullflexible,
}

\title{\textbf{Settlemint: a self-settling event market for NEAR}}
\author{Dipankar Sarkar \\ \small doom2quake}
\date{August 2026}

\begin{document}
\maketitle

\begin{abstract}
An on-chain event market that resolves a real-world question has a settlement
problem that predates crypto: someone decides the outcome, someone pays every
winner, and everyone else must be able to check that the payout matched the
outcome and that nothing was stranded. Settlemint is a NEAR-native contract that
makes settlement mechanical and provably complete. A resolver opens a binary
market, stakers commit yoctoNEAR while it is open, and when the outcome is known
the resolver records it once and every winner is paid pro rata in a single
resumable pass. The last winning claimant sweeps the rounding remainder, so
escrow lands at exactly zero and no stake is stranded, and a permissionless void
escape hatch lets every staker recover their own stake if the resolver never
resolves. This paper describes the settlement state machine, its conservation
argument, and the milestone-1 implementation: a NEAR Rust contract tested to 23
unit checks and compiled to a 205\,KB testnet WASM, and a Python mirror of the
same arithmetic held in cross-language parity by a further 32 tests. All work is
NEAR testnet only; no mainnet, no users, no audit.
\end{abstract}

\section{Introduction}

When a prediction market, an event contract, or an agent-run payout settles, a
participant who lost typically has no way to tell an honest resolution from a
rigged one, and no guarantee that every winner was paid to the last unit. Two
distinct failures hide inside ``the market settled.'' First, a \emph{liveness}
failure: if the resolver key goes silent, escrow can sit locked with no escape.
Second, an \emph{accounting} failure: pari-mutuel payout is a division, and naive
floor division drops a remainder, so the sum of payouts silently falls short of
the sum of stakes and some winner is quietly short-changed.

Settlemint targets both failures directly on NEAR. It is deliberately narrow: it
does not attempt to decide subjective outcomes (that is a later milestone), only
to make the \emph{settlement} mechanical, complete, and checkable once an outcome
is recorded. This paper covers the milestone-1 core: the on-chain state machine,
its conservation property, and the tested implementation.

\section{Why NEAR}

Two NEAR features make this design fit the chain rather than merely run on it.

\paragraph{Scoped access keys.} NEAR accounts separate full-access keys from
named function-call access keys~\cite{near_docs_access_keys_2026}. An operator
agent that settles markets can hold a key restricted to the \texttt{resolve} and
\texttt{settle} methods of the settlement contract, so a compromised agent can
advance markets but cannot move funds out of an account. The blast radius of a
lost key is exactly the two methods, not the whole balance.

\paragraph{Predictable low fees and yocto accounting.} NEAR settles in
yoctoNEAR ($10^{-24}$ NEAR) with predictable gas~\cite{near_nomicon_2026}, so the
one-resolve-plus-one-settle-pass pattern is cheap enough to run at cadence, and
the dust that a pari-mutuel split produces is a genuine integer remainder the
contract must account for rather than a floating-point artifact. The contract is
written in Rust against \texttt{near-sdk}~\cite{near_sdk_rs_2026} and compiled to
WASM; it is a reimplementation of the state machine for the NEAR runtime, not a
transliteration of EVM bytecode.

NEAR Protocol Rewards~\cite{near_protocol_rewards_2026} scores measured on-chain
activity (contract calls, unique wallets) directly, so an agent that settles many
small testnet markets produces exactly the honest, measured activity the
programme is built to reward.

\section{The settlement state machine}

An event is a record with a named resolver, a close timestamp, a phase, an
outcome, the two side pools, the escrow still held, and the winning-side stake
not yet claimed. It moves through four phases.

\begin{lstlisting}
Open      create_event names a resolver and closes_at; take_position adds yocto
Resolved  resolver records the outcome, only at/after close; winning pool frozen
Settled   every winner paid pro rata; reached only when escrow == 0
Voided    resolver silent past close + 7 days; each staker withdraws own stake
\end{lstlisting}

\paragraph{Resolver-only, post-close resolution.} Only the named resolver may
call \texttt{resolve}, and only at or after \texttt{closes\_at}. Resolving early
would let the resolver cut off an open market mid-flight, so the contract forbids
it.

\paragraph{Pro-rata payout with a dust sweep.} On a resolved market, an account
that staked $m$ on the winning side, where the winning pool is $W$ and the total
pool is $T$, is owed
\[
  \text{payout} = \left\lfloor \frac{T \cdot m}{W} \right\rfloor,
\]
except that the \emph{last} unclaimed winner receives the entire remaining
escrow instead of its floor share. Because escrow is decremented by every payout
and the final winner takes whatever is left, the rounding remainder is paid out
rather than stranded.

\paragraph{Refund cases.} If the winning side had no stake, or the market was
voided, each account is refunded exactly its own total stake. Voiding is
permissionless after \texttt{closes\_at + 7\,days}, which is why a lost or silent
resolver key cannot lock funds.

\section{Conservation}

Let $S$ be the sum of all stakes on an event. The contract maintains an
\texttt{escrow} field equal to the yoctoNEAR still held, initialised to $S$ as
positions are taken and decremented by each payout. Two invariants hold.

\begin{enumerate}
\item \textbf{No overpayment.} Each account is marked \texttt{claimed} before its
transfer is scheduled, so no account is paid twice, and each floor share is at
most its exact proportional share, so the running total of floor shares never
exceeds $S$.
\item \textbf{No underpayment.} The last unclaimed winner is paid the entire
remaining escrow, which is $S$ minus the sum of the earlier floor shares. Hence
the total paid equals $S$ exactly, and \texttt{escrow} reaches $0$.
\end{enumerate}

A market is reported \texttt{Settled} only once \texttt{escrow} reaches zero, so
a partly-paid market is never mistaken for a final one while a winner is still
owed money. The implementation asserts this with a test that settles one of two
winners and checks the phase is still \texttt{Resolved} with non-zero escrow.

\section{Implementation and evidence}

\paragraph{The contract.} \texttt{contract/src/lib.rs} is the whole mechanism in
one file, depending only on \texttt{near-sdk}. It compiles to a 205\,KB release
WASM with \texttt{cargo build --release --target wasm32-unknown-unknown}, which is
the artifact a testnet deploy uploads.

\paragraph{On-chain tests.} \texttt{cargo test} runs \textbf{23} unit tests on
the NEAR \texttt{unit-testing} harness, covering every phase transition, the
access-control checks, the pro-rata split, the dust sweep, the no-winner and void
refunds, and the not-finalised-until-escrow-zero property.

\paragraph{Off-chain mirror and cross-language parity.}
\texttt{settlemint/model.py} reimplements the identical arithmetic in Python.
\texttt{tests/test\_parity.py} asserts the exact payout literals the Rust tests
assert, so a one-line drift in either implementation fails a test. The Python
suites total \textbf{32} passing tests, including a run-fixture test that checks
the recorded run in \texttt{docs/run.json} conserves value to the last yocto and
reproduces exactly, and an adapter test that proves the live NEAR client refuses
any mainnet endpoint.

\begin{table}[h]
\centering
\begin{tabular}{lll}
\toprule
Suite & Command & Result \\
\midrule
NEAR Rust core & \texttt{cargo test} & 23 passing \\
Python model, parity, fixtures & \texttt{pytest tests -q} & 32 passing \\
Deployable artifact & \texttt{cargo build --release} & 205\,KB WASM \\
\bottomrule
\end{tabular}
\caption{Milestone-1 evidence, run 2026-08-26, all offline and keyless.}
\end{table}

\section{Limitations}

All work is NEAR testnet only; the live client rejects mainnet endpoints in code.
The recorded run is an offline fixture that executes the settlement arithmetic
but touches no network, so no public block-explorer receipt exists yet; this repo
ships no funded key. The verifiable-resolution hash commitment, the operator
agent on a scoped access key, and the forkable template are later milestones and
are not in this release. There are no users, no revenue, no partnerships, and no
third-party audit; the 55 passing tests are the authors' own. These limits are
restated in \texttt{docs/LIMITATIONS.md}.

\section{Related work}

Pari-mutuel and market-scoring-rule mechanisms have a long
history~\cite{hanson2003combinatorial}, and decentralized prediction markets such
as Augur~\cite{peterson2019augur} put resolution and payout on-chain. Settlemint
is narrower and more mechanical: it does not decide subjective outcomes in this
milestone, it makes the \emph{settlement} step complete and conservative on NEAR,
with a liveness escape hatch and an integer dust sweep as first-class properties.

\bibliographystyle{plain}
\bibliography{references}

\end{document}
